\section{Introduction}
\label{sec:intro}

\IEEEPARstart{5}{G} network slicing is one of the defining operational
ideas of modern mobile networks: a single physical radio access and
transport infrastructure is carved into logically isolated virtual
networks, each engineered for a service class with its own key
performance indicators --- ultra-reliable low-latency communications
(URLLC), enhanced mobile broadband (eMBB), and massive machine-type
communications (mMTC)~\cite{rost2017network,afolabi2018slicing}.
The economic promise of slicing rests on a deceptively simple
operational primitive: the operator periodically decides how much of a
shared resource pool (radio spectrum, transport bandwidth, compute) to
\emph{reserve} for each slice or sector over the next control interval.
Every downstream commitment --- a guaranteed bitrate for an industrial
URLLC flow, a busy-hour throughput target for a video eMBB cohort ---
is ultimately backed by that reservation decision. When the reservation
meets or exceeds realized demand, users are served and the SLA holds;
when it falls short, the slice violates its contract; whatever is
reserved but unused is idle capital that neighboring tenants cannot
use.

The difficulty is that the reservation must be committed \emph{before}
demand realizes. Mobile traffic is strongly periodic (diurnal and
weekly cycles), heavy-tailed at 15-minute granularity, and subject to
abrupt regime shifts --- mass events, flash crowds, application
launches, holidays --- that no historical model fully
anticipates~\cite{shafiq2012cellular,sciancalepore2017mobile}. A reservation
policy is therefore an exercise in decision-making under uncertainty,
and the natural language for that exercise is probabilistic
forecasting: predict not a single demand number but a distribution, and
size the reservation from a calibrated upper quantile of that
distribution. Yet the majority of deployed practice still sizes
reservations from static rules of thumb --- peak-plus-margin being the
most common --- precisely because the probabilistic alternative has
been perceived as operationally fragile.

This paper argues that the perception is outdated, and that the
remaining fragility has a different root: not the quality of deep
probabilistic forecasts, which is now adequate, but the
\emph{interface} between forecasts and reservations, and the
\emph{information hygiene} of the learning loop that adapts them. We
present UCRA, a closed-loop framework that makes both explicit. UCRA
comprises five stages: (1) a quantile LSTM forecaster that emits a
full predictive distribution of demand one hour ahead; (2) an
\emph{uncertainty-to-reservation transform} $\Phi$ that converts the
forecast distribution into a reservation target through one auditable
equation with a single interpretable risk parameter $\kappa$; (3) a
risk-adaptive allocation that distributes an aggregate reservation
across competing sectors or slices in proportion to demand pressure
and risk; (4) an update-and-release mechanism with EWMA smoothing and
release hysteresis that suppresses reservation churn; and (5) a
self-evolving learning stage that watches the reservation feedback
stream (violations, demand z-scores), detects distribution drift, and
fine-tunes the forecaster on a reservoir-sampled replay buffer ---
under an information-flow protocol that guarantees the loop never
learns from labels a live system could not yet possess.

\subsection{Motivation: two failure modes of the obvious designs}

The two canonical baselines fail in instructive, opposite directions,
and our experiments quantify both on real data. A \emph{static
peak-plus-margin} reservation --- the historical maximum demand, held
constant --- never violates on our three-week evaluation window, but
wastes 30.5\% of capacity; on a three-week trace that is roughly a
week of capacity paid for and never used. Worse, its safety is
conditional on stationarity: had a single test-window spike exceeded
the training peak, the ``safe'' policy would have violated with no
recourse. A \emph{median-forecast} reservation --- reserve what the
point forecast says demand will be --- is the opposite failure: it
tracks demand tightly (its mean absolute tracking error is the lowest
of all policies we test) yet violates the reservation in 48.0\% of
15-minute intervals on the same window, because the median of a
right-skewed distribution is by construction exceeded half the time.
The lesson is structural: violation rate and utilization trade off
along a frontier, and the frontier's shape is determined by how
uncertainty information enters the reservation rule. UCRA's
contribution is to make that entry point explicit, tunable, and
adaptive.

\subsection{Motivation: leakage is the silent killer of online-learning claims}

Self-evolving (continually learning) forecasters are increasingly
proposed for network operations, and our own Stage~5 is one. But
online-learning experiments carry a methodological trap that is easy
to fall into and easy to miss: at loop time $t$ a forecaster that
predicts an $H$-step horizon holds labels only for slots up to $t$;
its multi-step training targets for slots $t+1,\dots,t+H-1$ do not
\emph{exist} yet. An evaluation harness that stores full multi-step
windows into a replay buffer at time $t$ --- as ours originally did
--- lets the fine-tuned model peek as much as $H-1$ slots into the
future, inflating post-adaptation performance in a way that no live
deployment could realize. The pre-submission audit of this paper
caught exactly this flaw in our own pipeline; the final version
therefore introduces \emph{delayed causal replay}, which admits a
training window into the buffer only when its last label has been
observed, and refreshes forecasts only for slots whose reservations
have not yet been committed. We formalize the rule, prove its
sufficiency, unit-test it, and re-run every experiment under it. We
highlight this not as a detail but as a reusable evaluation protocol
for the growing literature on self-adaptive network
controllers~\cite{kapoor2023leakage}.

\subsection{A running example}

To make the framework concrete before the formal treatment, consider
one operational day on the evaluated network. At 05:45 the forecast
for the 06:00 slot reads a median of 61{,}200 with a wide
morning-ramp dispersion of 12{,}400; the transform reserves
$\approx$77{,}500 --- well above the median, because the forecaster
itself reports high uncertainty about how fast the morning ramp will
run. At 13:30, demand is stable and the band tightens; the same
formula reserves barely 6\% above the base quantile. At 21:00 a
flash crowd begins: reservations start violating, the rolling
violation indicator rises, the risk term inflates every buffer by up
to 45\% within its bounded range, and if the surge persists past the
drift trigger the forecaster is fine-tuned --- on strictly past
labels only --- so the next morning's forecasts reflect the new
regime. Every one of these decisions is a four-term arithmetic
expression an operator can audit after the fact, and the paper's
central empirical claim is that this loop, so specified, held zero
violations for three weeks of live traffic at two-thirds utilization
while static sizing wasted nearly a third of capacity.

\subsection{Contributions}

This paper makes the following contributions:

\begin{enumerate}
\item \textbf{An auditable uncertainty-to-reservation transform.} A
closed-form mapping $\Phi$ from a quantile forecast to a reservation,
$\Rt = \qtau + \kappa\,\spread\,(1 + w\,\rhot)$, with clamping to
a minimum-headroom floor and a capacity ceiling. One dial, $\kappa$,
moves the operator along the risk--utilization frontier; we trace the
full frontier empirically.
\item \textbf{A five-stage closed loop with self-evolution.} Quantile
LSTM forecasting, risk-feedback state, reservation smoothing with
release hysteresis, and trigger-based continual learning from a
reservoir replay buffer, wired into one control loop with per-slot
feedback.
\item \textbf{A formally causal online-learning protocol.} The
delayed replay intake (admission at $t+H-1$) and partial forecast
refresh (only slots $>t$), with a sufficiency argument, unit tests
that fail on the leaky variant, and a head-to-head quantification
showing the leaky variant inflates adaptation gains.
\item \textbf{A verified, multi-dataset empirical study.} On live RAN
performance counters: 0.0\% violations at 72.1\% utilization (seed
42; $67.6\%\pm3.4$ across three seeds) versus 30.5\% waste for static
peak sizing and 48.0\% violations for median-forecast sizing; a
cheating oracle still violates 16.3\%. Under a $+35\%$ surge: causal
self-evolution halves post-surge violations ($34.3\%\pm5.9
\rightarrow 17.2\%\pm4.9$). On the independent CTTC slicing workload:
1.2--2.7\% slice violations versus 5.8--43.7\% for the operator's
default sizing.
\item \textbf{An honest negative-control analysis.} The clean RAN
window triggers \emph{zero} self-evolution updates in all seeds: the
drift detector disciplines itself when there is nothing to fix. We
report this as a false-positive control and scope all online-learning
claims to the drift experiment, rather than over-claiming from the
clean run.
\item \textbf{Full reproducibility.} Deterministic configurations,
pinned datasets (Zenodo), unit tests including causality regression
tests, and one-command reruns of every table and figure in this
paper.
\end{enumerate}

\subsection{Paper organization}

Section~\ref{sec:related} reviews related work on slicing resource
management, predictive uncertainty, and online adaptation.
Section~\ref{sec:novelty} states the novelty claims and their
relationship to prior art. Section~\ref{sec:system} formalizes the
system model and the reservation problem. Section~\ref{sec:method}
details the five stages, including the causal replay protocol.
Section~\ref{sec:setup} describes datasets, baselines, metrics, and
the verification protocol. Section~\ref{sec:results} presents the
verified results. Section~\ref{sec:discussion} discusses scope,
limitations, and threats to validity, and Section~\ref{sec:conclusion}
concludes.
